% Figure — Amplificateur faible bruit à dégénérescence inductive (a) et schéma équivalent aux petites variations découpé en quadripôles A et β_F (b) (EANA3, chapitre 3).
% (a) D'après le support [support, p. 6, figure 1], redessiné.
% (b) D'après le corrigé externe [annale-2016-2017, p. 1], redessiné.
% Rendu SVG :
%   latex -output-directory=<dossier temporaire> 03-lna-degenerescence.tex
%   dvisvgm --no-fonts --exact-bbox -o 03-lna-degenerescence.svg <dossier temporaire>/03-lna-degenerescence.dvi
\documentclass[border=4pt]{standalone}
\usepackage[T1]{fontenc}
\usepackage[utf8]{inputenc}
\usepackage{amsmath}
\usepackage{sansmath}
\renewcommand{\familydefault}{\sfdefault}
\usepackage[european, straightvoltages, americancurrents]{circuitikz}
\begin{document}
\sansmath
\begin{circuitikz}[font=\small, >=latex, fleche/.style={->, thick}]
  % ---------- (a) ----------
  \node at (-0.4,4.6) {(a)};
  \draw (2.4,1.6) node[nmos] (M) {};
  \node[right] at (M.text) {$M$};
  \draw (M.D) -- ++(0,0.3) coordinate (d) to[R, l=$R_L$] ++(0,1.8) node[vcc] {$V_{DD}$};
  \draw (d) -- (4.2,0 |- d) node[ocirc]{};
  \draw (M.S) to[cute inductor, l=$L$] (M.S |- 0,-0.6) node[ground]{};
  \draw (0,0 |- M.G) node[ocirc]{} -- (M.G);
  \draw (0,-0.6) node[ocirc]{} node[ground]{};
  \draw[fleche] (-0.3,-0.45) -- node[left]{$V_{in}$} (-0.3,1.45);
  \draw[fleche] (4.5,-0.45) -- node[right]{$V_{out}$} (4.5,0 |- d);
  % ---------- (b) ----------
  \begin{scope}[xshift=7.0cm]
    \node at (-0.4,4.6) {(b)};
    \draw[thick, fill=black!3] (0.8,1.3) rectangle (5.0,4.1); \node[font=\bfseries] at (4.7,3.85) {$A$};
    \draw[thick, fill=black!3] (0.8,-1.3) rectangle (5.0,0.6); \node at (4.6,-1.05) {$\beta_F$};
    \draw (0,3.4) node[ocirc]{} -- (1.8,3.4) node[circ]{} node[above]{G} to[C, l=$C_{gs}$] (1.8,1.8) -- (3.4,1.8);
    \draw[fleche] (0.3,3.65) -- node[above]{$i_e$} (1.1,3.65);
    \draw[fleche] (1.25,1.95) -- node[left]{$v_{gs}$} (1.25,3.25);
    \draw (3.4,3.4) to[american controlled current source, l=$g_m v_{gs}$] (3.4,1.8);
    \draw (3.4,3.4) node[circ]{} node[above]{D} -- (6.2,3.4) node[ocirc]{};
    \draw[fleche] (6.0,3.65) -- node[above]{$i_o$} (5.3,3.65);
    \draw (2.6,1.8) node[circ]{} node[below right]{S} -- (2.6,0.3) to[cute inductor, l=$L$] (2.6,-1.7);
    \draw[fleche] (2.0,-1.5) -- node[left]{$v_{fb}$} (2.0,0.3);
    \draw[fleche] (3.1,0.2) -- node[right]{$i_L = i_o + i_e$} (3.1,-0.6);
    \draw (0,-1.7) node[ocirc]{} -- (6.2,-1.7) node[ocirc]{};
    \draw (3.6,-1.7) node[ground]{};
    \draw[fleche] (-0.3,-1.5) -- node[left]{$V_{in}$} (-0.3,3.2);
    \draw[fleche] (6.5,-1.5) -- node[right]{$V_{out}$} (6.5,3.2);
  \end{scope}
\end{circuitikz}
\end{document}
